\chapter{LITERATURE REVIEW}
\justifying

This chapter provides an extensive overview of existing research relevant to the three problem domains discussed in this work. The problem domains include: (i) log-based anomaly detection in distributed systems, (ii) secure image sharing and reconstruction using visual cryptography and learning-based methods, and (iii) cryptanalysis of chaos-based medical image encryption schemes. Although these problem domains address distinct problem contexts, they share common requirements with regard to learning representations from structured or semi-structured data and evaluating system security under advanced adversarial models. As such, this chapter is broadly divided into three sections, each covering one of these problem domains.
%--------------------------------------------------
\section{Anomaly Detection in Logs}

Log anomaly detection has been extensively studied as an important technique in guaranteeing reliability and security in large-scale distributed systems. Initially, researchers applied mining-based techniques to console logs using statistical and clustering-based approaches \cite{xu2009detecting, lou2010mining, makanju2009clustering}. However, such techniques heavily relied on structured log templates. To overcome this problem, researchers developed new techniques such as Drain \cite{he2017drain} and LogCluster \cite{vaarandi2015logcluster} to automatically parse structured log templates. In addition, researchers developed new tools to standardize parsing-based techniques \cite{zhu2019tools}. However, such techniques face challenges in terms of accuracy due to their dependency on parsing-based techniques.

To overcome this problem, learning-based techniques were developed to detect anomalies in logs. DeepLog \cite{du2017deeplog} employed an LSTM neural network to model sequential dependencies in logs to perform anomaly detection by learning to predict the next event in the logs. LogAnomaly \cite{meng2019loganomaly} further improved DeepLog to detect sequential and quantitative anomalies in logs. The problem of performing robust detection in unstable logs has been addressed in \cite{zhang2019robust}. More recently, researchers have proposed new parsing-free techniques such as Log2Vec \cite{meng2020log2vec} and parsing-free anomaly detection frameworks \cite{neurallog2021}. The transformer-based technique has been developed to model logs based on the attention mechanism \cite{vaswani2017attention} and BERT-based contextual embedding techniques \cite{devlin2018bert}. The increasing importance of attention-based and large language models in log analysis is demonstrated by recent works like LAnoBERT \cite{lee2023lanobert}, LogFormer \cite{guo2024logformer}, and several LLM-based parsing studies \cite{wu2023effectiveness, le2023prompt, ma2024llmparser, jiang2024lilac, xu2024divlog, xiao2024logbatcher}. Extensive surveys \cite{landauer2023survey} emphasize the necessity of robust and structurally aware log representation techniques.

\section{Attack on Secure Image Sharing Schemes using Autoencoders}

The concept of secret image sharing was first motivated by the development of classical secret sharing schemes by Shamir \cite{shamir1979share}, which laid a theoretical foundation for secret sharing among multiple parties. Naor and Shamir \cite{naor1995visual} later extended this idea to visual cryptography for images. The main idea is to encode a secret image into a number of shares such that any qualified set of shares will reveal the secret image visually. However, any unqualified subset will reveal nothing. The problem of visual cryptography is that there is a loss in contrast and resolution. 

To address this problem, XOR-based secret image sharing schemes have been proposed. The XOR-based secret image sharing schemes operate on image pixels instead of bits. The XOR-based secret image sharing schemes use Shannon's perfect secrecy concept \cite{shannon1949communication} to guarantee perfect reconstruction of the secret image from all the shares. The XOR-based secret image sharing schemes also guarantee that any subset of shares is statistically independent of the secret image. This ensures that there is zero mutual information between any unqualified subset of shares and the secret image.

However, the security proofs for the XOR-based constructions assume the presence of adversaries without access to structural priors about natural images. Recent deep learning advances contradict this. Convolutional Neural Networks (CNNs) \cite{lecun1998gradient} have shown impressive performance in learning hierarchical spatial representations from image data. Architectural innovations such as batch normalization \cite{ioffe2015batch} improve the stability of the learning process, while encoder-decoder architectures, such as the U-Net \cite{ronneberger2015unet}, enable accurate reconstruction for inverse image problems. Autoencoders, when learned using pixel-wise reconstruction objectives, can extract compact feature representations that can reconstruct structured image information from corrupted images. Moreover, the use of Perceptual Loss functions \cite{johnson2016perceptual} can improve the reconstruction quality for images. The reconstruction quality for such images can be measured using full-reference image similarity measures. The Peak Signal-to-Noise Ratio (PSNR) can measure the pixel-wise fidelity of the reconstruction, while the Structural Similarity Index Measure (SSIM) \cite{wang2004image} can measure the structural similarity between images by comparing the luminance, contrast, and structural information.

In spite of the solid theoretical foundation of secret sharing via XOR and the remarkable advancements in the field of image reconstruction via deep learning methods, only a few works have attempted to explore the possibility of reconstructing the secret image via learning the priors of images and utilizing only certain subsets of the shares generated via XOR operations. Most of the existing works have been dedicated to enhancing the efficiency and robustness of secret image sharing methods and verifying their feasibility rather than testing their robustness against learning-based attacks. This work attempts to bridge this gap and test the robustness of secret image sharing via XOR operations against learning-based attacks and determine whether the secret image can be reconstructed via learning the priors of images and utilizing only certain subsets of the shares generated via XOR operations.

\section{Cryptanalysis of Chaos-Based Image Encryption Schemes}

Chaos-based image encryption has been widely studied due to its ability to achieve strong confusion and diffusion properties using nonlinear dynamical systems. Early work by Fridrich \cite{fridrich1998image} introduced the permutation-diffusion architecture based on chaotic maps, which has since become a standard design paradigm. In this framework, permutation is used to break spatial correlation among pixels, while diffusion ensures that small changes in the plaintext propagate across the ciphertext.

Several schemes have been proposed by combining different chaotic systems. Guan et al. \cite{guan2005chaos} utilized Arnold's Cat Map along with Chen's chaotic system to enhance permutation complexity. Wu et al. \cite{wu2012image} explored logistic maps for both permutation and diffusion processes. Later works incorporated higher-dimensional chaotic systems and hybrid approaches to improve randomness and enlarge the key space \cite{mursi2014image, hua2018lscm, jain2021medical}. These schemes often report strong statistical performance based on metrics such as entropy, correlation coefficients, NPCR, and UACI.

Despite these improvements, a significant body of research has shown that many chaos-based encryption schemes suffer from fundamental security weaknesses. Xie et al. \cite{xie2017cryptanalysis} demonstrated vulnerabilities in permutation-diffusion structures due to insufficient coupling between stages. Cokal and Solak \cite{cokal2009cryptanalysis} showed that certain chaotic image ciphers are vulnerable to known-plaintext and chosen-plaintext attacks. More recent studies \cite{wen2023cryptanalysis, wen2024cryptanalysis, wen2023lin} have emphasized that plaintext-independent keystream generation is a major design flaw, as it allows attackers to recover equivalent keys using simple algebraic analysis.

A recurring issue in many schemes is that the keystream used in the diffusion stage depends only on the secret key and not on the plaintext. This results in a reusable keystream, which is equivalent to a one-time pad used multiple times. Such designs are inherently vulnerable to chosen-plaintext attacks, where specially crafted inputs can be used to extract the keystream and invert the encryption process. Additionally, the separability of permutation and diffusion stages allows attackers to independently recover pixel mappings and value transformations.

The scheme proposed by Jain et al. \cite{jain2021medical}, which combines Arnold's Cat Map with a Two-Dimensional Logistic Sine Coupling Map, represents a recent attempt to improve security for medical image encryption. However, similar to earlier works, it does not incorporate plaintext-dependent key generation or sufficient coupling between stages. This makes it susceptible to structural cryptanalysis.

The present work builds upon these observations and demonstrates a complete chosen-plaintext attack on such a scheme. By recovering the equivalent keystream using a zero plaintext and reconstructing the permutation mapping using structured probe images, the encryption can be fully inverted without knowledge of the secret key. This highlights the gap between statistical security evaluation and actual cryptographic robustness, reinforcing the need for rigorous cryptanalysis in the design of chaos-based image encryption schemes.
%--------------------------------------------------